% ===========================================================================
% Acknowledgments
% ===========================================================================
\section*{Acknowledgments}
\addcontentsline{toc}{section}{Acknowledgments}

This work was developed within the ``One-Person Company'' (一人公司)
framework, a 9-role AI-augmented software engineering methodology.
Each role contributed distinct expertise:

\begin{itemize}[nosep]
    \item \textbf{天明 (CTO)} — architecture vision and mathematical theory roadmap
    \item \textbf{心怡 (CPO)} — product requirements and agricultural domain knowledge
    \item \textbf{浩轩 (Staff Engineer)} — core implementation (7,329 lines, 313 tests)
    \item \textbf{子涵 (QA)} — quality metrics and topological validation
    \item \textbf{志远 (Data)} — preprocessing pipeline and spectral analysis
    \item \textbf{晓彤 (Code Review)} — adversarial verification and code quality
    \item \textbf{大鹏 (DevOps)} — deployment infrastructure and CI/CD
    \item \textbf{老王 (Storage)} — database schema and experiment tracking
    \item \textbf{雨桐 (UX)} — frontend design and user interaction
\end{itemize}

The mathematical formalization benefited from Lean 4's proof assistant for
formal verification of all 10 core theorems. The engineering implementation
was accelerated by Claude Code's multi-agent orchestration and the ALJ
Knowledge Engine for persistent context management.

The SIREN architecture is based on the work of Sitzmann et al. (NeurIPS 2020).
The EDL framework builds upon Sensoy et al. (NeurIPS 2018). Persistent
homology computation follows the algorithms of Edelsbrunner and Harer (2008).
Remote sensing pre-training uses the SeCo method of Ma\~{n}as et al. (2022).

This research was conducted at the School of Remote Sensing and Information
Engineering, Wuhan University. The authors acknowledge the Copernicus
programme for Sentinel-1/2 data and NASA for SRTM and Landsat data,
made freely available through Google Earth Engine and the AWS Open Data
programme.
